\chapter{Língua Inglesa}
\label{cap:ingles}

\begin{edital}
Compreensão de textos em inglês e itens gramaticais relevantes para o
entendimento dos sentidos dos textos.
\end{edital}
\peso{12 questões, peso 1 --- mesmo peso do Português. É quase toda
interpretação de texto. Para quem lê documentação técnica, é a disciplina de
melhor retorno por hora estudada de toda a prova.}

\section{O que a FGV cobra}

Foram 12 questões na Dataprev 2024 (Q13--Q24), todas ancoradas em
\textbf{textos reais} em inglês (anúncio de e-book, resenha de produto,
abstract acadêmico). Um texto serve a 2--6 questões. Os itens que mais
aparecem:

\begin{enumerate}
\item \textbf{Ideia principal / propósito do texto} (``what information is
in the text'', ``the main purpose is''). Você tem que captar a ideia central
e o gênero (é anúncio? resenha? artigo acadêmico?).
\item \textbf{Referência de pronome:} a quem/quê o \textit{it, this, they,
her} se refere --- volte à frase anterior.
\item \textbf{Conectivos / discourse markers:} o que eles sinalizam.
\item \textbf{Vocabulário em contexto:} sentido de uma palavra no texto;
palavra que funciona como verbo $\times$ substantivo (ex.: \textit{browse,
rate}).
\item \textbf{Verdadeiro/falso sobre o texto} (``only 1 and 4 are true'').
\end{enumerate}

\section{Conectivos que caem muito}

\begin{center}
\begin{tabular}{@{}p{5.4cm}p{5.4cm}@{}}
\toprule
\textbf{Marcador} & \textbf{Sentido} \\
\midrule
thus, therefore, hence, so & consequência/efeito \\
however, whereas, but, although, despite & contraste/concessão \\
moreover, furthermore, and, besides & adição \\
because, since, as & causa \\
unless & a menos que (condição negativa) \\
meanwhile & enquanto isso (tempo) \\
\bottomrule
\end{tabular}
\end{center}

Exemplo que já caiu: \textit{``Thus''} introduz efeito/consequência;
\textit{``and''} como adição $\approx$ \textit{moreover}.

\section{Modais e tempos (para o sentido)}

\textbf{should} = sugestão/conselho; \textbf{must} = obrigação;
\textbf{can/could} = possibilidade/permissão; \textbf{may/might} =
possibilidade. O foco não é decorar regra, mas entender o \textbf{sentido}
que o item dá ao texto.

\begin{pegadinha}
Formato ``julgue as afirmativas 1--4'' e escolha o conjunto verdadeiro/falso.
O distrator inverte um detalhe: ``It is a printed book'' quando o texto diz
``eBook Kindle only'', ou ``tips only for fresh graduates'' quando o texto
atende também a quem tem 20 anos de experiência.

A alternativa de compreensão global mais \textbf{longa e bem escrita}
costuma ser a certa, mas cuidado: a errada adiciona uma palavra que o texto
não sustenta (``abroad'', ``internship'', ``social status'', ``web
developers'') --- extrapolação clássica, igual à do Português.

No pronome, o distrator oferece o substantivo mais \textbf{próximo} ou mais
``óbvio'', não o correto. \textit{``her''} pode se referir a ``mom'', não a
students/professors só porque estão perto no texto.

No conector, a banca confunde causa com contraste. \textit{``Thus''} =
\textbf{effect} (consequência), não example nem condition. Decore os pares:
thus/hence/so (consequência), whereas/while (contraste),
moreover/furthermore (adição), unless (condição negativa).
\end{pegadinha}

\begin{comosair}
É prova de \textbf{LEITURA}, não de conversação. Não precisa entender cada
palavra; localize a informação que a questão pede e releia \textbf{só}
aquela frase. Para pronome/vocabulário, volte à linha exata do trecho
citado --- a resposta está no texto, nunca no seu palpite. Elimine
alternativas que trazem informação \textbf{não mencionada} (mesmo que
plausível no mundo real); se o texto não afirma, é falsa --- mesma técnica do
português. Monte uma minitabela mental de conectivos e de modais antes da
prova; são pontos quase gratuitos e caem sempre. Nas questões de V/F com
afirmativas numeradas, resolva \textbf{uma afirmativa por vez} e vá cortando
as combinações de resposta.
\end{comosair}

\begin{jacaiu}
\textbf{Em prova real da FGV:} ideia principal do texto (guia de carreira em
TI); referência do pronome \textit{her}; conector \textit{thus} (efeito);
\textit{and} $\approx$ \textit{moreover}; \textit{should} (sugestão);
\textit{browse}/\textit{rate} como verbo $\times$ substantivo; referência de
\textit{it} (impressora) --- \textbf{Dataprev 2024} (Q13--Q24). Vale notar:
a Dataprev 2024 é a \textbf{única prova do corpus com Língua Inglesa} --- MPU,
TJ-RJ e ALERO 2026 não têm o bloco. Toda a calibração de estilo do capítulo
vem, portanto, dessas 12 questões.

\textbf{No nosso banco} (previsto pelo edital, ainda não visto na amostra de
provas): textos técnicos sobre nuvem, IA na programação, trabalho remoto,
\emph{phishing} e IA no setor público, com o mesmo desenho --- ideia
principal, referência, conectivo e vocabulário em contexto. Os dois lotes
mais recentes partem de \textbf{texto real}, no mesmo perfil de fonte que a
Dataprev 2024 usou: um \emph{abstract} do \textbf{arXiv} sobre o efeito medido
da IA na produtividade de desenvolvedores (\emph{forecast} $\times$
\emph{estimate}, \emph{Surprisingly} como contraexpectativa, \emph{This
slowdown} como anáfora) e um post do \textbf{Cisco Talos} sobre uso de
ferramenta de IA em campanha de \emph{phishing} contra a administração
pública (\emph{Notably} como realce, \emph{This incident} como anáfora,
\emph{lower the barrier to entry} como vocabulário em contexto). A fonte
completa de cada texto está no \texttt{why} da primeira questão do bloco.

Rode \texttt{./quiz.py ingles}.
\end{jacaiu}

\section{Pontos de atenção / pesquisa extra}

\begin{itemize}
\item \textbf{Vocabulário técnico de TI} ajuda muito (deploy, framework,
request, query) --- você já tem vantagem por estudar a área.
\item \textbf{Skimming} (ideia geral) e \textbf{scanning} (achar informação
específica) são as técnicas de leitura para ganhar tempo.
\end{itemize}
